\section{Manutenção predial sustentável no ambiente construído}
\label{sec:revisao}

O tratamento da manutenção predial como estratégia de sustentabilidade, e não apenas como despesa operacional, decorre de sua posição no ciclo de vida do ambiente construído: intervenções preventivas e preditivas prolongam a vida útil dos ativos, reduzem o consumo de recursos associado a substituições prematuras e mitigam riscos operacionais que, sem gestão adequada, tendem a se acumular como passivo técnico e financeiro. \textcite{ojobo_estrategiassustentaveis_2024}, em revisão de escopo internacional, associam estratégias de manutenção sustentável ao aumento da longevidade das edificações, reforçando a leitura de que a manutenção não é evento isolado, mas processo contínuo com efeito direto sobre indicadores ambientais, econômicos e sociais do ativo construído.

A busca conduzida para este artigo identificou, além do núcleo analítico principal, um corpus complementar tecnológico (registros associados a \textit{smart campus}, BIM, \textit{digital twin}, IoT e manutenção preditiva) e um corpus conceitual exploratório (registros associados a ambiente construído como sistema vivo, metabolismo urbano e biofilia), ambos deliberadamente mantidos fora do núcleo analítico principal. Esses dois blocos subsidiam contextualização e identificação de lacunas, mas não compõem a base empírica dos resultados apresentados nas seções~\ref{sec:panorama} a~\ref{sec:aplicabilidade}; qualquer menção a eles, ao longo do texto, é explicitamente qualificada como apoio contextual, não como achado do corpus principal.
